\documentclass[10pt]{beamer}

\usetheme{metropolis}
\definecolor{WiLabRed}{RGB}{197,18,48}
\setbeamercolor{frametitle}{fg=white,bg=WiLabRed}
\setbeamercolor{progress bar}{fg=WiLabRed!90}
\setbeamercolor{title separator}{fg=WiLabRed!90}
\setbeamercolor{progress bar in section page}{fg=WiLabRed!90}
\setbeamercolor{background canvas}{bg=white}
\setbeamercolor{alerted text}{fg=WiLabRed!90}

\usepackage{appendixnumberbeamer}

\usepackage{booktabs}
\usepackage[scale=2]{ccicons}

\usepackage{pgfplots}
\usepgfplotslibrary{dateplot}

\usepackage{xspace}
\newcommand{\themename}{\textbf{\textsc{metropolis}}\xspace}

\usepackage{marvosym}
%\usepackage{subfig}
\usepackage{graphicx}\graphicspath{{images/}}
\usepackage{subcaption}
\usepackage[framed]{./support/mcode}
\usepackage{listings}

\title{Lecture 04: Gaussian Noise Channels}
\subtitle{\textit{Software Defined Radio for Engineers} (Collins~\textit{et al.}), \textsection{3.5}}
\date{}
\author{\textbf{Alexander M. Wyglinski, Ph.D.}}
\institute{ \vspace*{1in}\hfill\includegraphics[height=1.125cm]{wilab_logo-A70916.eps} \qquad \includegraphics[height=1.125cm]{WPI_Inst_Prim_FulClr.eps}}
% \titlegraphic{\hfill\includegraphics[height=1.5cm]{logo.pdf}}

% Foot for all slides
\setbeamertemplate{frame footer}{\tiny \copyright~2018 by Alexander Wyglinski. This work is licensed under the Creative Commons Attribution-ShareAlike 4.0 International License. To view a copy of this license, visit http://creativecommons.org/licenses/by-sa/4.0/.}

\begin{document}

%\captionsetup[subfigure]{labelformat=empty}

%%%%%%%%%%%%%%%%%%%%%%%%%%%%%%%%%%%%%%%%%%%%%%%%%%%%%%%%%%

\maketitle




%%%%%%%%%%%%%%%%%%%%%%%%%%%%%%%%%%%%%%%%%%%%%%%%%%%%%%%%%%

\frame
{
  \frametitle{Noise Channel Model}
  
  \centering
  Binary Information Source $\rightarrow$ $m[n]$\\
  $\downarrow$\\
  Transmitter $\rightarrow$ $s(t)$\\
  $\downarrow$\\
  Channel $\rightarrow$ $s(t)+n(t)$\\
  $\downarrow$\\
  Receiver $\leftarrow$ $r(t)$\\
  $\downarrow$\\
  Binary Information Sink $\leftarrow$ $\hat{m}[n]$

}




%%%%%%%%%%%%%%%%%%%%%%%%%%%%%%%%%%%%%%%%%%%%%%%%%%%%%%%%%%

\frame
{
  \frametitle{Gaussian Random Variable}
  
  \begin{itemize}
   \item We frequently use \textbf{Gaussian random variables} to model noise contribution to transmitted signal
   \item Gaussian random variable mathematically expressed as:
   \begin{equation}
    f(x)=\frac{1}{\sqrt{2\pi{\sigma^2}}}e^{-((x-\mu)/\sigma)^2/2}
   \end{equation}
   where $\mu$ is the mean and $\sigma$ is the standard deviation
  \end{itemize}

  
}




%%%%%%%%%%%%%%%%%%%%%%%%%%%%%%%%%%%%%%%%%%%%%%%%%%%%%%%%%%

\frame
{
  \frametitle{Bivariate Gaussian}

General definition for bivariate
Gaussian density with parameters $\mu_X$, $\mu_Y$, $\sigma_X^2$,
$\sigma_Y^2$, and correlation coefficient $\rho$ is given by:
\begin{equation}
f_{XY}(x,y)=\frac{\exp\left(\frac{-1}{2(1-\rho^2)}\left(\left(\frac{x-\mu_X}{\sigma_X}\right)^2-2\rho\left(\frac{x-\mu_X}{\sigma_X}\right)\left(\frac{y-\mu_Y}{\sigma_Y}\right)+\left(\frac{y-\mu_Y}{\sigma_Y}\right)^2\right)\right)}{2\pi\sigma_X\sigma_Y\sqrt{1-\rho^2}},
\end{equation}
where {\it correlation coefficient} is defined as:
\begin{equation}
    \rho=E\left[\left(\frac{x-\mu_X}{\sigma_X}\right)\left(\frac{y-\mu_Y}{\sigma_Y}\right)\right].
\end{equation}
  
  
}
  

%%%%%%%%%%%%%%%%%%%%%%%%%%%%%%%%%%%%%%%%%%%%%%%%%%%%%%%%%%

\frame
{
  \frametitle{Zero Correlation Scenario}
  
  \centering
  What happens when $\rho=0$?\\
  $\downarrow$\\
  $f_{XY}(x,y)=f_{X}(x)f_{Y}(y)$\\
  \vspace*{1cm}
  $\rightarrow$ How does this look like?
  
  
  }
  
  

%%%%%%%%%%%%%%%%%%%%%%%%%%%%%%%%%%%%%%%%%%%%%%%%%%%%%%%%%%

\begin{frame}[fragile]{Impact of Gaussian Noise Channels}
  
\begin{figure}
\centering
\begin{minipage}[framed]{0.9\textwidth}
\begin{lstlisting}
% Noise standard deviation
nstddev = 0.1;

% Create data stream of information
wavefm = [];
for ind=0:1:9,
  if(round(rand(1,1))==1)
    wavefm = [wavefm ones(1,1000)];
  else
    wavefm = [wavefm zeros(1,1000)];  
  end;
end;

% Add noise to the transmission
noisysig = wavefm + nstddev.*randn(1,length(wavefm));
\end{lstlisting}
\end{minipage}
\captionsetup{labelformat=empty}
\end{figure}

\end{frame}

%%%%%%%%%%%%%%%%%%%%%%%%%%%%%%%%%%%%%%%%%%%%%%%%%%%%%%%%%%







\end{document}
